% META: {"id":"1SPE-SUITES-CO-017","chapitre":"1SPE-SUITES","type_objet":"corrige","exercice_ref":"1SPE-SUITES-EX-017","capacites":["1SPE-SUITES-C2","1SPE-SUITES-C6"],"capacites_codes":["C2","C6"],"sources_inspiration":[],"status":"approved","origin":"generated","status_history":["generated","audited","accepted","approved"],"release_acceptance":"RELEASE_OWNER_BATCH_ACCEPTANCE","acceptance_closure_digest":"sha256:9b3ccf9a81c5520fbb7e03b7b2d3e7bf2057a02834b6d908bf3be5f49f3b553f","acceptance_receipt_digest":"sha256:4fad86f0282e0fa4e0419cca1ad6379ae7af5a280c04a4a90880f90a1c044242"}
% BEGIN-VERIFY
% from sympy import *
% n = symbols('n', integer=True, nonnegative=True)
% C = 500 + 30*n
% assert C.subs(n, 0) == 500
% assert C.subs(n, 1) == 530
% assert C.subs(n, 12) == 860
% # Suite arithmetique de raison 30
% diff_C = C.subs(n, n+1) - C
% assert diff_C.simplify() == 30
% # Premier terme C0 = 500, raison r = 30
% END-VERIFY
\begin{corrige}{1SPE-SUITES-EX-017}

\commentaireMarge{On modélise la situation : le coût total est la somme du versement initial et des mensualités accumulées.}

\textbf{Question 1.} Après $n$ mois, Yasmine a payé $500$~€ à l'inscription et $30 \times n$~€ de mensualités. Donc :
\[
C_n = 500 + 30n \quad \text{(en euros).}
\]

\commentaireMarge{Une suite est arithmétique si et seulement si la différence de deux termes consécutifs est constante.}

\textbf{Question 2.} On calcule $C_{n+1} - C_n$ :
\[
C_{n+1} - C_n = \bigl(500 + 30(n+1)\bigr) - (500 + 30n) = 30.
\]

La différence $C_{n+1} - C_n = 30$ est constante pour tout $n \in \mathbb{N}$, donc la suite $(C_n)$ est \textbf{arithmétique} de :
\begin{itemize}
  \item premier terme $C_0 = 500 + 30 \times 0 = 500$~€ ;
  \item raison $r = 30$.
\end{itemize}

\commentaireMarge{On applique la formule du terme général de la suite arithmétique $C_n = C_0 + nr$ avec $n = 12$.}

\textbf{Question 3.} On calcule le coût total après $12$ mois :
\[
C_{12} = 500 + 30 \times 12 = 500 + 360 = 860.
\]

Après $1$ an ($12$ mois), le coût total est de $\mathbf{860}$~\textbf{€}.

\end{corrige}
